\chapter{Nouns}

\section{Inflection}

\subsection{Number}

Interturkic nouns distinguish singular from plural morphologically.
The plural is marked with the suffix \foreign{\latntxt{\=/lAr}}, though it is
not added if plurality is already indicated by other means, such as numerals or
adverbs of quantity; compare (\ref{ex:plural-lar-yes}) and
(\ref{ex:plural-lar-no}).

\pex
  \a\label{ex:plural-lar-yes}
    \vtop{\halign{
      \foreign{\latntxt{#}}\hfil&\:\transl{#}\hfil&\:\derivedto~\foreign{\latntxt
      {#}}\hfil\cr
      köz&eye&közler\cr
      kitab&book&kitablar\cr
    }}
  \a\label{ex:plural-lar-no}
    \vtop{\halign{
      \foreign{\latntxt{#}}\hfil&\:\transl{#}\hfil\cr
      eki köz&two eyes\cr
      tört kitab&four books\cr
    }}
\xe

By contrast, the singular is marked by the absence of the plural suffix.
In addition to denoting actual singularity, a noun in the singular can also have
generic or ambiguous readings, as in (\ref{ex:generic-singular}).
Singularity can be expressed more explicitly by adding the indefinite article
\foreign{\latntxt{bir}} \transl{a(n)}, as in (\ref{ex:explicit-singular}).

\pex
  \a\label{ex:generic-singular}%
    \begingl
      \gla at\=/\zero{} kütmek//
      \glb horse-\Sg{} {to herd}//
      \glft\transl{to herd a horse} or \transl{to herd horses}//
    \endgl
  \a\label{ex:explicit-singular}%
    \begingl
      \gla Bir kitab\=/\zero{} oquyman.//
      \glb one book-\Sg{} {I read}//
      \glft\transl{I read a book.}//
    \endgl
\xe

\subsection{Possession}
\label{sec:possession}

Possession of nouns is expressed by the suffixes in table~\ref{%
tbl:posessive-suffixes}, which encode the person and number of the possessor.
Note that the third\-/person suffix is the same for both singular and plural
possessors.

\begin{table}
  \centering
  \caption{Possessive suffixes}
  \label{tbl:posessive-suffixes}
  \begin{tabular}{llll}
    \toprule
    &\First&\Second&\Third\\
    \midrule
    \Sg&\latntxt{\=/(I)m}&\latntxt{\=/(I)ñ}
    &\multirow[t]{2}{*}{\latntxt{\=/(s)I}}\\
    \cmidrule(lr){2-3}
    \Pl&\latntxt{\=/(I)m}&\latntxt{\=/(I)ñIz}\\
    \bottomrule
  \end{tabular}
\end{table}

\subsection{Case}

Case describes a noun's role within a phrase and is encoded by the suffixes in
table~\ref{tbl:case-suffixes}.
Note that the nominative is encoded by the lack of a case suffix, and that
certain case suffixes have different forms when preceded by possessive suffixes.

\begin{table}
  \centering
  \caption{Case suffixes}
  \label{tbl:case-suffixes}
  \begin{tabular}{llll}
    \toprule
    &\multirow[t]{2}{*}{Default}&\multicolumn{2}{l}{Possessive}\\
    \cmidrule(lr){3-4}
    &&\First/\Second&\Third\\
    \midrule
    \Nom&\multicolumn{3}{l}{\latntxt{\=/\zero}}\\
    \Acc&\multicolumn{2}{l|}{\latntxt{\=/nI}}&\latntxt{\=/n(I)}\\
    \Gen&\multicolumn{3}{l}{\latntxt{\=/nIñ}}\\
    \Dat&\latntxt{\=/GA}&\multicolumn{2}{|l}{\latntxt{\=/(n)A}}\\
    \Loc&\latntxt{\=/DA}&\multicolumn{2}{|l}{\latntxt{\=/(n)dA}}\\
    \Abl&\latntxt{\=/DAn}&\multicolumn{2}{|l}{\latntxt{\=/(n)dAn}}\\
    \bottomrule
  \end{tabular}
\end{table}

The nominative marks subjects~(\ref{ex:nom-subj}), vocatives~(\ref{ex:nom-voc}),
copular complements~(\ref{ex:nom-cop-compl}), and indefinite direct objects~(%
\ref{ex:nom-indef-dir-obj}).

\pex
  \a\label{ex:nom-subj}%
    \begingl
      \gla Ahmad\=/\zero{} kitab oquydu.//
      \glb Ahmad\=/\Nom{} book reads//
      \glft\transl{Ahmad reads a book.}//
    \endgl
  \a\label{ex:nom-voc}%
    \begingl
      \gla Ey watandaşlar\=/\zero{}, qaydasız?//
      \glb \Voc{} compatriot.\Pl\=/\Nom{}, where.\Cop.\Spl{}//
      \glft\transl{O compatriots, where are you?}//
    \endgl
  \a\label{ex:nom-cop-compl}%
    \begingl
      \gla Sen Iranlı\=/\zero\=/san.//
      \glb \Ssg{} Iranian\=/\Nom\=/\Cop.\Ssg{}//
      \glft\transl{You're Iranian.}//
    \endgl
  \a\label{ex:nom-indef-dir-obj}%
    \begingl
      \gla Çatır\=/\zero{} al!//
      \glb umbrella\=/\Nom{} take//
      \glft\transl{Take an umbrella!}//
    \endgl
\xe

The accusative marks direct objects if and only if they are considered to be
definite (see \ref{sec:definiteness}).

\ex
  \begingl
    \gla Kitab\=/ı\=/n oqu\=/ma\=/dı\=/ñ\=/mı?//
    \glb book\=/\Poss.\Third\=/\Acc{} read\=/\Neg\=/\Pst\=/\Ssg\=/\Q//
    \glft\transl{Haven't you read his book?}//
  \endgl
\xe

The genitive marks possessors~(\ref{ex:gen-poss}) and the complements of certain
postpositions~(\ref{ex:gen-compl}).
The usage in (\ref{ex:gen-poss}) mirrored by the possessive suffixes added to
the possessum (see \ref{sec:possession}).

\pex
  \a\label{ex:gen-poss}%
    \begingl
      \gla meniñ yüzüg\=/üm//
      \glb \Fsg.\Gen{} ring\=/\Poss.\Fsg{}//
      \glft\transl{my ring}//
    \endgl
  \a\label{ex:gen-compl}%
    \begingl
      \gla üy\=/nüñ iç\=/i\=/nde//
      \glb house\=/\Gen{} interior\=/\Poss.\Third\=/\Loc//
      \glft\transl{inside the house}//
    \endgl
\xe

The dative marks indirect objects~(\ref{ex:dat-indir-obj}), places towards which
an action is directed~(\ref{ex:dat-dest}), the complements of certain
postpositions~(\ref{ex:dat-post-compl}), and the complements of certain verbs~%
(\ref{ex:dat-verb-compl}).

\pex
  \a\label{ex:dat-indir-obj}%
    \begingl
      \gla Oña suw berdim.//
      \glb \Tsg.\Dat{} water {I gave}//
      \glft\transl{I gave him water.}//
    \endgl
  \a\label{ex:dat-dest}%
    \begingl
      \gla Yapraq yer\=/ge tüştü.//
      \glb leaf ground\=/\Dat{} fell//
      \glft\transl{The leaf fell to the ground.}//
    \endgl
  \a\label{ex:dat-post-compl}%
    \begingl
      \gla yel\=/ge qarşı//
      \glb wind\=/\Dat{} against//
      \glft\transl{against the wind}//
    \endgl
  \a\label{ex:dat-verb-compl}%
    \begingl
      \gla Maña bağıñız.//
      \glb \Fsg.\Dat{} look//
      \glft\transl{Look at me.}//
    \endgl
\xe

The locative expresses location within some space, whether that be a physical~%
(\ref{ex:loc-phys}), temporal~(\ref{ex:loc-time}), or some abstract concept~(%
\ref{ex:loc-abstr}).

\pex
  \a\label{ex:loc-phys}%
    \begingl
      \gla Xacıtarxan\=/da yaşaylar.//
      \glb Astrakhan\=/\Loc{} {they live}//
      \glft\transl{They live in Astrakhan.}//
    \endgl
  \a\label{ex:loc-time}%
    \begingl
      \gla Hüseyn sağat beş\=/te keldi.//
      \glb Husayn hour five\=/\Loc{} left//
      \glft\transl{Husayn left at five o'clock.}//
    \endgl
  \a\label{ex:loc-abstr}%
    \begingl
      \gla Tınçlıq\=/ta edik.//
      \glb peace\=/\Loc{} {we were}//
      \glft\transl{We were in [a time of] peace.}//
    \endgl
\xe

The ablative marks points of departure~(\ref{ex:abl-origin}) and places
`through' which an action is performed~(\ref{ex:abl-through}), as well the
complements of certain verbs~(\ref{ex:abl-verb-compl}) and postpositions~(\ref{%
ex:abl-post-compl}).

\pex
  \a\label{ex:abl-origin}%
    \begingl
      \gla Tawlar\=/dan keldiler.//
      \glb mountain.\Pl\=/\Abl{} {they came}//
      \glft\transl{They came from the mountains.}//
    \endgl
  \a\label{ex:abl-through}%
    \begingl
      \gla Ağaçnı pencere\=/den kördüm.//
      \glb tree.\Acc{} window\=/\Abl{} {I saw}//
      \glft\transl{I saw the tree through the window.}//
    \endgl
  \a\label{ex:abl-verb-compl}%
    \begingl
      \gla Ondan qorqma.//
      \glb \Tsg.\Abl{} {be not afraid}//
      \glft\transl{Don't be afraid of it.}//
    \endgl
  \a\label{ex:abl-post-compl}%
    \begingl
      \gla ders\=/ten soñ//
      \glb lession\=/\Abl{} after//
      \glft\transl{after the lesson}//
    \endgl
\xe

\subsection{Definiteness}
\label{sec:definiteness}

Unlike the other grammatical categories, definiteness is not encoded directly on
a noun.
Rather, a noun's context may confer definiteness upon it.
Nouns are considered definite when they are possessed, proper, or qualified by
a demonstrative.
